% REQUIRED CSSI SUPPLEMENT -- 2-page limit.
\section*{Delivery Mechanism and Community Usage Metrics}

\subsection*{Delivery Mechanisms}

\textbf{Hosted service.} CryoStack will be delivered through the Georgia Tech
gateway at \url{https://cryostack.eas.gatech.edu}. The gateway will provide
documentation, authenticated workspaces, saved configurations, experiment tracking, the
Frozen Legacies historical-radar catalog, LIVIST
(Living Ice Sheet Temperature) exploration using radar- and borehole-derived
products, CryoLauncher model execution, and ICESEE data assimilation.
Core laptop-scale tutorials will remain usable without remote HPC.

\textbf{Open software and environments.} Source, issues, releases, adapter
templates, and governance will be public in the ICESEE-project GitHub
organization. Twice-yearly tagged releases will be archived with DOI records.
Supported execution environments will be delivered as source recipes, Spack
locks, and OCI/Apptainer images pinned by digest, with software bills of
materials, license manifests, compatibility matrices, and validation reports.

\textbf{Remote execution.} Typed, allow-listed execution adapters will stage
declared inputs, submit and monitor jobs, and recover declared outputs on
qualified backends. Georgia Tech PACE is the reference backend; AWS Batch (Fargate by default,
EC2 On-Demand as an opt-in profile) provides a portable cloud profile for
project-managed or user-managed accounts, allocations, and credits; CryoStack
does not require AWS, and outside budgeted project activities AWS bills the
researcher's own account. A bounded access pilot, if approved, gives up to 50
participants a small introductory allocation for eligible CryoLauncher and
ICESEE workflows. The application-level experiment specification
will remain consistent across profiles. Credentials remain with the user
connector or an approved credential service, not in the public gateway
database.

\textbf{Community delivery.} The project will offer one annual hands-on
workshop, quarterly office hours in Years 2--3, recorded/literate tutorials,
and a public support channel. A contribution guide, adapter proposal template,
code of conduct, roadmap, support policy, and governance process will make the
service usable and extensible beyond the originating team.

\subsection*{Yearly Targets and Collection}

Targets are evaluation commitments, not claims of present adoption. ``Unique
users'' are privacy-preserving counts of authenticated people who complete at
least one scientific run in the reporting year. ``External group'' means a
non-project research or teaching group with at least two completed runs or a
documented curricular/research use. Automated traffic and project-team tests
are excluded from adoption counts.

\begin{table}[ht]
\centering
\scriptsize
\begin{tabular}{@{}p{0.32\textwidth}rrrp{0.25\textwidth}@{}}
\toprule
Metric & Y1 & Y2 & Y3 & Collection/evidence \\
\midrule
Qualified model adapters (cumulative) & 4 & 5 & 5 & CryoBench registry and release tags; PISM is the fifth \\
Qualified dataset adapters (cumulative) & 1 & 2 & 2 & Validated manifests, provenance/QC reports, and release tags \\
Adapters passing required release gates & 80\% & 100\% & 100\% & Signed CI/benchmark reports \\
Qualified remote backend profiles (cumulative) & 1 & 2 & 2 & PACE and AWS qualification checklists \\
PISM-specific changes confined to the adapter/configuration layer (shared-code changes only via documented extension interfaces) & -- & $\geq$90\% & $\geq$90\% & Code-review diff classification at production release \\
Tagged stack releases per year & 2 & 2 & 2 & GitHub/Zenodo release records \\
Core-workflow availability, excluding planned maintenance & 95\% & 97\% & 97\% & External uptime monitor (Year-1 instrument; current checks are internal) and incident log \\
Median time to acknowledge high-severity service issues & 2 d & 1 d & 1 d & Public/private issue timestamps \\
Annual unique scientific users & 25 & 60 & 100 & De-identified account/run records \\
Completed non-test scientific runs per year & 150 & 500 & 1,000 & Experiment registry; bot/team tests excluded \\
External adopting groups (cumulative) & 2 & 5 & 8 & Opt-in affiliation plus qualifying usage/interview \\
Annual workshop participants & 25 & 40 & 50 & Registration and attendance records \\
Core user-task completion rate, each of the four applications & 65\% & 75\% & 80\% & Application-specific task studies \\
Access-pilot first-run completion (if pilot approved) & -- & 70\% & 80\% & Pilot run records; cost per completed workflow \\
Allocation enforcement tests passed (attribution, refusal, pool isolation, recovery) & -- & 100\% & 100\% & Task 2.4 test reports; reconciled estimate-vs-actual cost \\
Median System Usability Scale score & 65 & 70 & 75 & Anonymous post-study survey \\
External contributors (cumulative) & 1 & 2 & 3 & Accepted code, docs, test, or adapter contributions \\
Research/teaching outputs citing CryoStack (cumulative) & 1 & 3 & 6 & User survey, DOI/citation alerts, manual verification \\
\bottomrule
\end{tabular}
\end{table}

\subsection*{Quality, Science, and Broader-Impact Evidence}

Each production adapter must pass schema/API tests; deterministic restart and
seeded reproducibility tests; a model-specific scientific benchmark; an
end-to-end assimilation test with predefined error/spread tolerances; and a
cross-backend scientific-equivalence test on every qualified backend it
claims to support, comparing field norms, integrated or conserved
quantities, ensemble statistics, or DA-specific diagnostics against
versioned tolerance envelopes rather than requiring matching runtime or cost.
The PISM adapter must complete a synthetic-truth experiment and a documented
observational demonstration. Benchmark reports will record code/environment
versions, runtime, resource use, and tolerances.

Workshop reporting will include attendance, career stage, institution type,
completion of the hands-on workflow, and an anonymous usefulness survey.
Demographic information will be optional, minimized, reported only in
aggregate, and suppressed for small cells. Accessibility checks and user-task
studies will generate issue lists and closure evidence. Annual public reports
will compare targets with actuals, explain shortfalls, identify corrective
actions, and publish machine-readable aggregates where privacy permits.

\subsection*{Instrumentation and Responsibility}

The durable identity/experiment registry and expanded Control Center are the
sources for users, runs, backends, status, audit events, availability, and
reproducibility fields; they will store no scientific credentials.
GitHub and Zenodo APIs provide release/contribution evidence. An external
uptime monitor and incident log provide operations data. Registration forms,
task-study instruments, and short annual opt-in user surveys provide community
evidence. The PI reviews technical and operations metrics monthly; the Co-PI
reviews scientific and user metrics quarterly. Frozen yearly extracts,
definitions, exclusion rules, and analysis scripts will be archived with each
annual report so that trends remain auditable if instrumentation changes.
